\scanpage{11}

Réciproquement, partons d'une équation comme dessus, où les \deleted{$a$}\inserted{coeff. $a,b\ldots$} sont des sections de $\underline O_S$. On sait qu'\deleted{$a$}\inserted{une telle équation} définit un diviseur relatif \inserted{positif} de $P_S^2/S$, soit $C$. Il reste à écrire que c'est là une \uline{courbe simple} sur $S$. On sait qu'il faut pour cela que le discriminant de la forme
\[
\delta(a,b,c,d,e,f)
\]
soit une \uline{unité} partout. Ici le discriminant s'explicite : ---

On a obtenu en résumé

\uline{Proposition} La donnée d'une courbe elliptique sur $S$, \deleted{\unsure{avec}}\inserted{munie} d'une \deleted{``base adaptée''}\inserted{``base spéciale''} de $\mathcal E_3(C/S)$, \deleted{``unité''} est équivalente à la donnée de sections $a,b,c,d,e,f$ de $\underline O_S$ satisfaisant $\delta(a,b,c,d,e,f)\unsure{\neq|=}0$ ; cette correspondance est compatible avec l'image inverse.

On tire de ceci : Considérons le schéma affine de type fini sur $\mathbb Z$
\[
M=\operatorname{Spec}(\mathbb Z[A,B,C,D,E,F][\delta^{-1}])
\]
qui est un ouvert affine de $\operatorname{Spec}\mathbb Z[A,B,C,D,E,F]\simeq \unsure{\mathbb E|E}_{\mathbb Z}^6$. \deleted{Faisons opérer \unsure{dessus|dans}}
\mspage{10}
